\secciongrande{Capacidades y limitaciones de ejecución ¿por qué javascript en el cliente}

\resumen{
  El desarrollo web dinámico se gestiona por lenguajes de script en el lado del servidor, pero JavaScript es esencial en el cliente para mejorar la usabilidad. Sus capacidades incluyen la preparación de datos, la navegación funcional sin JS y la actualización dinámica de contenido recibido del servidor.
}

\definicion{JavaScript en el cliente}{Lenguaje de script ejecutado en el navegador del usuario, enfocado en la interacción y mejora de la experiencia de usuario (UX) sin acceder al sistema de archivos del servidor.} 
\subtituloA{Razones para usar JavaScript en el cliente}

\lista{
  \textbf{Preparación de datos:} JavaScript no puede escribir archivos en el entorno servidor; su función es preparar datos para su envío o delegar su tratamiento a otro lenguaje del servidor. |
  \textbf{Navegabilidad sin javascript:} El sitio debe ser funcional sin JS; este se usa para optimizar la navegación (rápida, moderna) sin dejar la página inactiva si JS está deshabilitado. |
  \textbf{Actualización dinámica:} Permite examinar datos recibidos del lado del servidor y actualizar selectivamente secciones de la página web sin recarga completa.
}
